\documentclass[11pt]{article}
\usepackage[T1]{fontenc}
\usepackage[utf8]{inputenc}
\usepackage[margin=1in]{geometry}
\usepackage{amsmath,amssymb,amsthm}
\usepackage{enumerate}

\newtheorem{theorem}{Theorem}[section]
\newtheorem{lemma}[theorem]{Lemma}
\newtheorem{corollary}[theorem]{Corollary}
\theoremstyle{definition}
\newtheorem{definition}[theorem]{Definition}
\newtheorem{example}[theorem]{Example}
\theoremstyle{remark}
\newtheorem*{remark}{Remark}

\newcommand{\N}{\mathbb{N}}
\newcommand{\Z}{\mathbb{Z}}
\newcommand{\R}{\mathbb{R}}
\DeclareMathOperator{\lcm}{lcm}

\title{Lecture 4: Divisibility and the Euclidean Algorithm}
\author{Course 21-127 --- Concepts of Mathematics}
\date{Week 2}

\begin{document}
\maketitle
\tableofcontents

\section{Divisibility}

\begin{definition}[Divides]
Let $a, b \in \Z$. We say that $a$ \emph{divides} $b$, written $a \mid b$,
if there exists $k \in \Z$ such that $b = ka$. If no such $k$ exists we
write $a \nmid b$.
\end{definition}

\begin{example}
$3 \mid 12$ since $12 = 4 \cdot 3$, but $5 \nmid 12$. Every integer divides
$0$, and $1$ divides every integer.
\end{example}

\begin{lemma}\label{lem:transitive}
If $a \mid b$ and $b \mid c$ then $a \mid c$.
\end{lemma}

\begin{proof}
Write $b = ka$ and $c = \ell b$ for some $k, \ell \in \Z$. Then
$c = \ell b = \ell(ka) = (\ell k) a$, and $\ell k \in \Z$, so $a \mid c$.
\end{proof}

\begin{lemma}\label{lem:linear}
If $a \mid b$ and $a \mid c$ then $a \mid (mb + nc)$ for all $m, n \in \Z$.
\end{lemma}

\begin{proof}
Write $b = ka$ and $c = \ell a$. Then $mb + nc = mka + n\ell a = (mk + n\ell)a$.
\end{proof}

\begin{remark}
Lemma~\ref{lem:linear} is the fact that makes the Euclidean algorithm
work: the set of common divisors of $a$ and $b$ is unchanged when $b$ is
replaced by $b - qa$.
\end{remark}

\section{The division theorem}

\begin{theorem}[Division with remainder]\label{thm:division}
Let $a \in \Z$ and $b \in \Z_{>0}$. There exist unique $q, r \in \Z$ with
\[
  a = qb + r \quad\text{and}\quad 0 \le r < b.
\]
\end{theorem}

\begin{proof}
\emph{Existence.} Consider the set
$S = \{\, a - kb : k \in \Z,\ a - kb \ge 0 \,\}$. The set is non-empty
(take $k = -|a|$), so by well-ordering it has a least element $r = a - qb$.
If $r \ge b$ then $r - b = a - (q+1)b \in S$ is smaller than $r$, a
contradiction. Hence $0 \le r < b$.

\emph{Uniqueness.} Suppose $a = qb + r = q'b + r'$ with $0 \le r, r' < b$.
Then $(q - q')b = r' - r$ and $|r' - r| < b$, so $b \mid (r' - r)$ forces
$r' - r = 0$, and therefore $q = q'$.
\end{proof}

\begin{corollary}
Every integer is either even or odd, and not both.
\end{corollary}

\section{Greatest common divisors}

\begin{definition}
For $a, b \in \Z$ not both zero, $\gcd(a, b)$ is the largest $d \in \Z_{>0}$
with $d \mid a$ and $d \mid b$. We also write $\lcm(a,b)$ for the least
positive common multiple.
\end{definition}

\begin{theorem}[Euclidean algorithm]
Let $a \ge b > 0$ and write $a = qb + r$ as in Theorem~\ref{thm:division}.
Then $\gcd(a, b) = \gcd(b, r)$.
\end{theorem}

\begin{proof}
By Lemma~\ref{lem:linear}, every common divisor of $a$ and $b$ divides
$r = a - qb$, and every common divisor of $b$ and $r$ divides $a = qb + r$.
The two pairs therefore have the same common divisors, hence the same
greatest one.
\end{proof}

\begin{example}
To compute $\gcd(252, 198)$:
\begin{align*}
  252 &= 1 \cdot 198 + 54, \\
  198 &= 3 \cdot 54 + 36, \\
   54 &= 1 \cdot 36 + 18, \\
   36 &= 2 \cdot 18 + 0.
\end{align*}
The last non-zero remainder is $18$, so $\gcd(252, 198) = 18$.
\end{example}

\begin{theorem}[B\'ezout]
For $a, b \in \Z$ not both zero there exist $m, n \in \Z$ with
$ma + nb = \gcd(a, b)$.
\end{theorem}

\begin{proof}[Proof sketch]
Run the Euclidean algorithm and substitute backwards; each remainder is an
integer combination of the previous two, hence of $a$ and $b$.
\end{proof}

\subsection*{Exercises}
\begin{enumerate}[(i)]
  \item Prove that $\gcd(n, n+1) = 1$ for every $n \in \N$.
  \item Find $m, n \in \Z$ with $252m + 198n = 18$.
  \item Show that $\gcd(a,b) \cdot \lcm(a,b) = ab$ for $a, b \in \Z_{>0}$.
\end{enumerate}

\end{document}
